% GENERATED FILE. Edit canonical lessons, then run npm run generate:books.
% canonical-source-hash: fnv1a64:07e2fa5caaeb2f2f
% canonical-lessons: MR-C346-balagne, MR-C346-visad-karne, MR-C346-khatri, MR-C346-padtalne, MR-C346-phetalne

\chapter{Checking Claims}
\label{ch:mr-checking-claims}
\hlchaptermodality{346}

\begin{chapteropening}
\textbf{By the end of this chapter:} \emph{I can say how sure I am and check what I am told.}

Complete the last lesson of chapter 346: I can say how sure I am and check what I am told.
\end{chapteropening}

\section[\texorpdfstring{\mr{बाळगणे} (bāḷagṇe)}{बाळगणे (bāḷagṇe)}]{\mr{बाळगणे} (bāḷagṇe) — to entertain (a hope); to keep}

\label{lesson:MR-C346-balagne}



\begin{quote}
Before the new one: say the Marathi for a summary, then the Marathi for in short.
\end{quote}

\subsection*{You'll want to know: \mr{बाळगणे}}
\textbf{\mr{बाळगणे}} — \emph{bāḷagṇe} — \textquotedblleft{}to entertain (a hope); to keep\textquotedblright{}.

\begin{quote}
\mr{त्याने} \mr{निकालाबद्दल} \mr{मोठी} \mr{आशा} \mr{बाळगली} \mr{होती}.
\end{quote}

\begin{quote}
\emph{tyāne nikālābaddal moṭhī āśā bāḷaglī hotī.}
\end{quote}

\begin{quote}
He had cherished high hopes about the result.
\end{quote}

Past, with -\mr{ने} (he …): \mr{त्याने} \mr{आशा} \mr{बाळगली}, \emph{tyāne āśā bāḷaglī}.

\mr{आशा} \mr{बाळगणे}, to cherish a hope; \mr{राग} \mr{बाळगणे}, to bear a grudge.

\begin{cousinweb}
\textbf{\mr{बाळगणे}} (\emph{bāḷagṇe}) is a Marathi verb whose origin is not settled.
\end{cousinweb}

\begin{grammarlens}[title={What you've built}]
One more word for giving reasons and opinions.
\end{grammarlens}

\subsection*{Your turn}
\begin{itemize}
\raggedright
  \item \emph{Say it:} \emph{bāḷagṇe}
  \item \emph{Say it:} \emph{bāḷagṇe}, once more
  \item \emph{Say it:} \emph{tyāne nikālābaddal moṭhī āśā bāḷaglī hotī.} — He had cherished high hopes about the result.
  \item \emph{Recall:} say the Marathi for a summary, then the Marathi for in short, then say \emph{bāḷagṇe} again
\end{itemize}

\begin{tcolorbox}[breakable,colback=teal!4,colframe=teal!35!black,title={Before you move on}]
What does \mr{बाळगणे} mean? (\textquotedblleft{}To entertain (a hope); to keep\textquotedblright{}.) Say it once more.
\end{tcolorbox}
\section[\texorpdfstring{\mr{विशद} \mr{करणे} (viśad karṇe)}{विशद करणे (viśad karṇe)}]{\mr{विशद} \mr{करणे} (viśad karṇe) — to elucidate}

\label{lesson:MR-C346-visad-karne}



\begin{quote}
Before the new one: say the Marathi for in short, then the Marathi for to entertain.
\end{quote}

\subsection*{You'll want to know: \mr{विशद} \mr{करणे}}
\textbf{\mr{विशद} \mr{करणे}} — \emph{viśad karṇe} — \textquotedblleft{}to elucidate\textquotedblright{}.

\begin{quote}
\mr{शिक्षकांनी} \mr{उदाहरण} \mr{देऊन} \mr{नियम} \mr{विशद} \mr{केला}.
\end{quote}

\begin{quote}
\emph{śikṣakānnī udāhraṇ deūn niyam viśad kelā.}
\end{quote}

\begin{quote}
The teacher explained the rule in detail with an example.
\end{quote}

Past, with -\mr{ने} (he …): \mr{त्याने} \mr{विशद} \mr{केलं}, \emph{tyāne viśad kelaṁ}.

Formal; in speech \mr{समजावून} \mr{सांगणे}.

\begin{cousinweb}
\textbf{\mr{विशद} \mr{करणे}} (\emph{viśad karṇe}) joins the tatsama \mr{विशद}, Sanskrit \emph{viśada}, \textquotedblleft{}clear, bright\textquotedblright{} (Monier-Williams), to \mr{करणे}, \textquotedblleft{}to make\textquotedblright{}: to make clear.
\end{cousinweb}

\begin{grammarlens}[title={What you've built}]
One more word for giving reasons and opinions.
\end{grammarlens}

\subsection*{Your turn}
\begin{itemize}
\raggedright
  \item \emph{Say it:} \emph{viśad karṇe}
  \item \emph{Say it:} \emph{viśad karṇe}, once more
  \item \emph{Say it:} \emph{śikṣakānnī udāhraṇ deūn niyam viśad kelā.} — The teacher explained the rule in detail with an example.
  \item \emph{Recall:} say the Marathi for in short, then the Marathi for to entertain, then say \emph{viśad karṇe} again
\end{itemize}

\begin{tcolorbox}[breakable,colback=teal!4,colframe=teal!35!black,title={Before you move on}]
What does \mr{विशद} \mr{करणे} mean? (\textquotedblleft{}To elucidate\textquotedblright{}.) Say it once more.
\end{tcolorbox}
\section[\texorpdfstring{\mr{खात्री} (khātrī)}{खात्री (khātrī)}]{\mr{खात्री} (khātrī) — certainty; assurance}

\label{lesson:MR-C346-khatri}



\begin{quote}
Before the new one: say the Marathi for to entertain, then the Marathi for to elucidate.
\end{quote}

\subsection*{You'll want to know: \mr{खात्री}}
\textbf{\mr{खात्री}} — \emph{khātrī} — \textquotedblleft{}certainty; assurance\textquotedblright{}.

\begin{quote}
\mr{तो} \mr{येईल} \mr{याची} \mr{मला} \mr{खात्री} \mr{आहे}.
\end{quote}

\begin{quote}
\emph{to yeīl yācī malā khātrī āhe.}
\end{quote}

\begin{quote}
I am sure that he will come.
\end{quote}

Feminine.

X \mr{ची} \mr{खात्री} \mr{असणे}, to be sure of X; \mr{खात्री} \mr{करणे}, to make sure.

\begin{cousinweb}
\textbf{\mr{खात्री}} (\emph{khātrī}) is a loan, usually traced to Persian \emph{khāṭir}, \textquotedblleft{}mind\textquotedblright{}, itself from Arabic: a mind set at rest.
\end{cousinweb}

\begin{grammarlens}[title={What you've built}]
One more word for giving reasons and opinions.
\end{grammarlens}

\subsection*{Your turn}
\begin{itemize}
\raggedright
  \item \emph{Say it:} \emph{khātrī}
  \item \emph{Say it:} \emph{khātrī}, once more
  \item \emph{Say it:} \emph{to yeīl yācī malā khātrī āhe.} — I am sure that he will come.
  \item \emph{Recall:} say the Marathi for to entertain, then the Marathi for to elucidate, then say \emph{khātrī} again
\end{itemize}

\begin{tcolorbox}[breakable,colback=teal!4,colframe=teal!35!black,title={Before you move on}]
What does \mr{खात्री} mean? (\textquotedblleft{}Certainty; assurance\textquotedblright{}.) Say it once more.
\end{tcolorbox}
\section[\texorpdfstring{\mr{पडताळणे} (paḍtāḷṇe)}{पडताळणे (paḍtāḷṇe)}]{\mr{पडताळणे} (paḍtāḷṇe) — to verify; to cross-check}

\label{lesson:MR-C346-padtalne}



\begin{quote}
Before the new one: say the Marathi for to elucidate, then the Marathi for certainty.
\end{quote}

\subsection*{You'll want to know: \mr{पडताळणे}}
\textbf{\mr{पडताळणे}} — \emph{paḍtāḷṇe} — \textquotedblleft{}to verify; to cross-check\textquotedblright{}.

\begin{quote}
\mr{पत्रकाराने} \mr{बातमी} \mr{छापण्याआधी} \mr{पडताळली}.
\end{quote}

\begin{quote}
\emph{patrakārāne bātmī chāpaṇyāādhī paḍtāḷlī.}
\end{quote}

\begin{quote}
The journalist checked the story before printing it.
\end{quote}

Past, with -\mr{ने} (he …): \mr{त्याने} \mr{पडताळलं}, \emph{tyāne paḍtāḷlaṁ}.

Transitive: the past agrees with the thing checked, \mr{बातमी} \mr{पडताळली}; \mr{पडताळा}, a cross-check.

\begin{cousinweb}
\textbf{\mr{पडताळणे}} (\emph{paḍtāḷṇe}) is a Marathi verb whose origin is not settled; the noun \mr{पडताळा}, a cross-check, goes with it.
\end{cousinweb}

\begin{grammarlens}[title={What you've built}]
One more word for giving reasons and opinions.
\end{grammarlens}

\subsection*{Your turn}
\begin{itemize}
\raggedright
  \item \emph{Say it:} \emph{paḍtāḷṇe}
  \item \emph{Say it:} \emph{paḍtāḷṇe}, once more
  \item \emph{Say it:} \emph{patrakārāne bātmī chāpaṇyāādhī paḍtāḷlī.} — The journalist checked the story before printing it.
  \item \emph{Recall:} say the Marathi for to elucidate, then the Marathi for certainty, then say \emph{paḍtāḷṇe} again
\end{itemize}

\begin{tcolorbox}[breakable,colback=teal!4,colframe=teal!35!black,title={Before you move on}]
What does \mr{पडताळणे} mean? (\textquotedblleft{}To verify; to cross-check\textquotedblright{}.) Say it once more.
\end{tcolorbox}
\section[\texorpdfstring{\mr{फेटाळणे} (pheṭāḷṇe)}{फेटाळणे (pheṭāḷṇe)}]{\mr{फेटाळणे} (pheṭāḷṇe) — to reject (a claim)}

\label{lesson:MR-C346-phetalne}



\begin{quote}
Before the new one: say the Marathi for certainty, then the Marathi for to verify.
\end{quote}

\subsection*{You'll want to know: \mr{फेटाळणे}}
\textbf{\mr{फेटाळणे}} — \emph{pheṭāḷṇe} — \textquotedblleft{}to reject (a claim)\textquotedblright{}.

\begin{quote}
\mr{कंपनीने} \mr{कामगारांची} \mr{मागणी} \mr{फेटाळली}.
\end{quote}

\begin{quote}
\emph{kampnīne kāmgārāñcī māgṇī pheṭāḷlī.}
\end{quote}

\begin{quote}
The company rejected the workers' demand.
\end{quote}

Past, with -\mr{ने} (he …): \mr{त्याने} \mr{मागणी} \mr{फेटाळली}, \emph{tyāne māgṇī pheṭāḷlī}.

Of claims, demands and charges; of a person or an offer, \mr{नाकारणे}.

\begin{cousinweb}
\textbf{\mr{फेटाळणे}} (\emph{pheṭāḷṇe}) is a Marathi verb whose origin is not settled.
\end{cousinweb}

\begin{grammarlens}[title={What you've built}]
One more word for giving reasons and opinions.
\end{grammarlens}

\subsection*{Your turn}
\begin{itemize}
\raggedright
  \item \emph{Say it:} \emph{pheṭāḷṇe}
  \item \emph{Say it:} \emph{pheṭāḷṇe}, once more
  \item \emph{Say it:} \emph{kampnīne kāmgārāñcī māgṇī pheṭāḷlī.} — The company rejected the workers' demand.
  \item \emph{Recall:} say the Marathi for certainty, then the Marathi for to verify, then say \emph{pheṭāḷṇe} again
\end{itemize}

\begin{tcolorbox}[breakable,colback=teal!4,colframe=teal!35!black,title={Before you move on}]
What does \mr{फेटाळणे} mean? (\textquotedblleft{}To reject (a claim)\textquotedblright{}.) Say it once more.
\end{tcolorbox}
